% ===========================================================================
% 2 Literature Review (about 9 pages): the 2022-2026 survey, 40-50 references.
% Drafted with AI assistance (Claude). Rewrite in your own words before submitting.
% ===========================================================================
\section{Literature Review}
\label{sec:literature}

\subsection{Lexical retrieval}
% TODO: BM25, probabilistic relevance framework; learned sparse (SPLADE).

\subsection{Dense retrieval}
% TODO: bi-encoders (Sentence-BERT, DPR), contrastive pre-training (Contriever),
% late interaction (ColBERT), general embedding models (BGE, M3), biomedical (MedCPT).

\subsection{Domain shift and benchmarks}
% TODO: BEIR, SciFact, SciFact-Open, LitSearch, BRIGHT; domain adaptation (GPL).

\subsection{Hybrid retrieval and rank fusion}
% TODO: RRF, analysis of fusion functions, zero-shot hybrid, the weakest-link effect
% (PVLDB 2026), dynamic weights (DAT), small embeddings with reranking.

\subsection{Neural reranking}
% TODO: monoBERT, monoT5, LLM rerankers, distilled cross-encoders, cost.

\subsection{Retrieval-augmented generation for the scientific literature}
% TODO: RAG, Fusion-in-Decoder, citations (ALCE), Self-RAG, CRAG, OpenScholar, PaperQA2.

\subsection{Literature assistants in practice}
% TODO: SciSpace, Consensus, ResearchRabbit: what they do, what they do not publish.

\subsection{Knowing when not to answer}
% TODO: selective QA, abstention survey, MultiVerS.

\subsection{Comparative summary}
% TODO: table of the surveyed approaches (task, data, strength, limitation).

\subsection{Research gap}
\label{sec:gap}
% TODO: aggregate scores only; no claim-level explanation of which retriever wins; no
% analysis of what a fusion weight does; reranking cost rarely weighed against
% in-domain fine-tuning; tools closed. This paper addresses each (Section~\ref{sec:novelty}).
